% Options for packages loaded elsewhere
\PassOptionsToPackage{unicode}{hyperref}
\PassOptionsToPackage{hyphens}{url}
%
\documentclass[
  a4paper,
  DIV=11,
  numbers=noendperiod]{scrartcl}

\usepackage{amsmath,amssymb}
\usepackage{iftex}
\ifPDFTeX
  \usepackage[T1]{fontenc}
  \usepackage[utf8]{inputenc}
  \usepackage{textcomp} % provide euro and other symbols
\else % if luatex or xetex
  \usepackage{unicode-math}
  \defaultfontfeatures{Scale=MatchLowercase}
  \defaultfontfeatures[\rmfamily]{Ligatures=TeX,Scale=1}
\fi
\usepackage{lmodern}
\ifPDFTeX\else  
    % xetex/luatex font selection
\fi
% Use upquote if available, for straight quotes in verbatim environments
\IfFileExists{upquote.sty}{\usepackage{upquote}}{}
\IfFileExists{microtype.sty}{% use microtype if available
  \usepackage[]{microtype}
  \UseMicrotypeSet[protrusion]{basicmath} % disable protrusion for tt fonts
}{}
\makeatletter
\@ifundefined{KOMAClassName}{% if non-KOMA class
  \IfFileExists{parskip.sty}{%
    \usepackage{parskip}
  }{% else
    \setlength{\parindent}{0pt}
    \setlength{\parskip}{6pt plus 2pt minus 1pt}}
}{% if KOMA class
  \KOMAoptions{parskip=half}}
\makeatother
\usepackage{xcolor}
\usepackage[textheight=9in,textwidth=6.5in,top=1in,headheight=14pt,headsep=25pt,footskip=30pt]{geometry}
\usepackage{svg}
\setlength{\emergencystretch}{3em} % prevent overfull lines
\setcounter{secnumdepth}{3}
% Make \paragraph and \subparagraph free-standing
\makeatletter
\ifx\paragraph\undefined\else
  \let\oldparagraph\paragraph
  \renewcommand{\paragraph}{
    \@ifstar
      \xxxParagraphStar
      \xxxParagraphNoStar
  }
  \newcommand{\xxxParagraphStar}[1]{\oldparagraph*{#1}\mbox{}}
  \newcommand{\xxxParagraphNoStar}[1]{\oldparagraph{#1}\mbox{}}
\fi
\ifx\subparagraph\undefined\else
  \let\oldsubparagraph\subparagraph
  \renewcommand{\subparagraph}{
    \@ifstar
      \xxxSubParagraphStar
      \xxxSubParagraphNoStar
  }
  \newcommand{\xxxSubParagraphStar}[1]{\oldsubparagraph*{#1}\mbox{}}
  \newcommand{\xxxSubParagraphNoStar}[1]{\oldsubparagraph{#1}\mbox{}}
\fi
\makeatother


\providecommand{\tightlist}{%
  \setlength{\itemsep}{0pt}\setlength{\parskip}{0pt}}\usepackage{longtable,booktabs,array}
\usepackage{calc} % for calculating minipage widths
% Correct order of tables after \paragraph or \subparagraph
\usepackage{etoolbox}
\makeatletter
\patchcmd\longtable{\par}{\if@noskipsec\mbox{}\fi\par}{}{}
\makeatother
% Allow footnotes in longtable head/foot
\IfFileExists{footnotehyper.sty}{\usepackage{footnotehyper}}{\usepackage{footnote}}
\makesavenoteenv{longtable}
\usepackage{graphicx}
\makeatletter
\newsavebox\pandoc@box
\newcommand*\pandocbounded[1]{% scales image to fit in text height/width
  \sbox\pandoc@box{#1}%
  \Gscale@div\@tempa{\textheight}{\dimexpr\ht\pandoc@box+\dp\pandoc@box\relax}%
  \Gscale@div\@tempb{\linewidth}{\wd\pandoc@box}%
  \ifdim\@tempb\p@<\@tempa\p@\let\@tempa\@tempb\fi% select the smaller of both
  \ifdim\@tempa\p@<\p@\scalebox{\@tempa}{\usebox\pandoc@box}%
  \else\usebox{\pandoc@box}%
  \fi%
}
% Set default figure placement to htbp
\def\fps@figure{htbp}
\makeatother

%% Palatino font:
\usepackage[T1]{fontenc}
\usepackage[utf8]{inputenc}
\usepackage{palatino}

\usepackage{pifont}

\usepackage{upquote}

\usepackage{fontspec}
\newfontfamily\mytitlefont{Lucida Grande}

% front page
\usepackage[colorlinks]{hyperref} % for text color on title
\usepackage{titling} % makes \thetitle etc available...
\usepackage{lipsum}
\usepackage{tikz}
\tikzstyle{bag} = [align=center] % allows new lines in text below
\AtBeginDocument{\thispagestyle{empty}
\definecolor{titlecolor}{HTML}{1C2404}
\begin{tikzpicture}[remember picture,overlay]
\draw (current page.center) node[inner sep=0, opacity=0.4] {\resizebox{!}{350mm}{\includegraphics{highres_banner.pdf}}};
\draw (current page.center) node[inner sep=0, opacity=0.5] {\resizebox{!}{100mm}{\includegraphics{ausegl_sort.pdf}}};
% \draw (current page.center) node[bag] {  
        \draw (17, 1.7) node[bag, align=right, anchor=north east] {  
    \mytitlefont \color{titlecolor} \Large Master of Science in Bioinformatics \\ 
    \\ 
    \mytitlefont \color{titlecolor} \Huge \thetitle \\ 
    \\ 
    \mytitlefont \color{titlecolor} \huge \theauthor \\
    \\
    \mytitlefont \color{titlecolor} \Large \thedate 
    };
\end{tikzpicture}
\clearpage}

% back side
\AtEndDocument{\clearpage\thispagestyle{empty}
\begin{tikzpicture}[remember picture,overlay]
\draw (current page.center) node[inner sep=0, opacity=0.4] {\resizebox{!}{350mm}{\includegraphics{highres_banner.pdf}}};
% \draw (current page.center) node[bag] {  
    \draw (17, 1.7) node[bag, align=right, anchor=north east] {  
        \mytitlefont \color{titlecolor} \large Bioinformatics Research Centre \\ 
        \mytitlefont \color{titlecolor} \large Department of Molecular Biology and Genetics \\ 
        \mytitlefont \color{titlecolor} \large Aarhus University \\
        \mytitlefont \color{titlecolor} \large Universitetsbyen 81 \\
        \mytitlefont \color{titlecolor} \large 8000 Aarhus C \\
        \mytitlefont \color{titlecolor} \large Denmark
        };

\end{tikzpicture}}


% picture (AU logo in title page)
\usepackage{titlepic}
\usepackage{graphicx}
\titlepic{\resizebox{!}{50mm}{\includegraphics[width=\textwidth]{ausegl_sort.pdf}}}



% \pagestyle{plain} % no running header

\usepackage[many]{tcolorbox}
\definecolor{quotegray}{HTML}{505050}
\newtcolorbox{myquote}{%
    enhanced jigsaw, 
    breakable,      % allow page breaks
    frame hidden,   % hide the default frame
    left=1cm,       % left margin
    right=1cm,      % right margin
    colback=white,
    fontupper=\color{quotegray},
    overlay={%
        \node [scale=4,
            text=lightgray,
            inner sep=0pt,] at ([xshift=0.5cm,yshift=-0.7cm]frame.north west){``}; 
        \node [scale=4,
            text=lightgray,
            inner sep=0pt,] at ([xshift=-0.5cm]frame.south east){''};  
            },
        % paragraph skips obeyed within tcolorbox
                parbox=false,
}
% redefine the 'quote' environment to use this 'myquote' environment
\renewenvironment{quote}{\begin{myquote}}{\end{myquote}}


% \usepackage{framed} % not sure i need this anymore

% \usepackage[T1]{fontenc}
% \usepackage{inconsolata}

% % I have to only define Shaded if it is already defined.
% % The reason is that pandoc does not define the macro if there are not code blocks in the a markdown file.
% \ifx \@Shaded \@empty

% \renewcommand{\KeywordTok}[1]{\textcolor[rgb]{0, 0, 0}{\textbf{{#1}}}} % def and or not reg
% \renewcommand{\BuiltInTok}[1]{\textcolor[rgb]{0.373, 0.298, 0.580}{\textbf{{#1}}}} % print open 
% \renewcommand{\VariableTok}[1]{\textcolor[rgb]{0.141, 0.392, 0.824}{\textbf{{#1}}}}
% \renewcommand{\OperatorTok}[1]{\textcolor[rgb]{0, 0, 0}{{#1}}} % def and or not reg
% \renewcommand{\DataTypeTok}[1]{\textcolor[rgb]{1.0,0.13,0.00}{{#1}}}
% \renewcommand{\DecValTok}[1]{\textcolor[rgb]{0.655, 0.498, 0.161}{{#1}}}
% \renewcommand{\BaseNTok}[1]{\textcolor[rgb]{0.259, 0.592, 0.596}{{#1}}}
% \renewcommand{\FloatTok}[1]{\textcolor[rgb]{0.655, 0.498, 0.161}{{#1}}}
% \renewcommand{\CharTok}[1]{\textcolor[rgb]{0.678,0.141,0.098}{{#1}}}
% \renewcommand{\StringTok}[1]{\textcolor[rgb]{0.678,0.141,0.098}{{#1}}}
% \renewcommand{\CommentTok}[1]{\textcolor[rgb]{0.135, 0.134, 0.133}{{#1}}}
% \renewcommand{\OtherTok}[1]{\textcolor[rgb]{0.00,0.44,0.13}{{#1}}}
% \renewcommand{\AlertTok}[1]{\textcolor[rgb]{1.00,0.00,0.00}{\textbf{{#1}}}}
% \renewcommand{\FunctionTok}[1]{\textcolor[rgb]{0.549, 0.102, 0.063}{\textbf{{#1}}}}  % function name
% \renewcommand{\RegionMarkerTok}[1]{{#1}}
% \renewcommand{\ErrorTok}[1]{\textcolor[rgb]{1.00,0.00,0.00}{\textbf{{#1}}}}
% \renewcommand{\NormalTok}[1]{\textcolor[rgb]{0, 0, 0}{{#1}}}


% \else
%   % no code blocks with markup...
% \fi



\usepackage{etoolbox}
\makeatletter
\g@addto@macro{\appendix}{%
  \patchcmd{\@@makechapterhead}% <cmd>
    {\endgraf\nobreak\vskip.5\baselineskip}% <search>
    {\hspace*{-.5em}:\space}% <replace>
    {}{}% <success><failure>
  \patchcmd{\@chapter}% <cmd>
    {\addchaptertocentry{\thechapter}}% <search>
    {\addchaptertocentry{Appendix~\thechapter:}}% <replace>
    {}{}% <success><failure>
  \addtocontents{toc}{%
    \protect\patchcmd{\protect\l@chapter}% <cmd>
      {1.5em}% <search>
      {6.5em}% <replace>
      {}{}}% <success><failure>
}
\renewcommand{\autodot}{}% Remove all end-of-counter dots
\makeatother


% restart chapter numbers in each part
\makeatletter
\@addtoreset{chapter}{part}
\makeatother


% \usepackage{chngcntr}
% \counterwithin*{subsubsection}{chapter}
% \counterwithout*{subsubsection}{section}
% \counterwithout*{subsubsection}{subsection}

% %  % KMT only use subsubsection number (this increments though the book 
% %  % when we supress numbering with  {.unnumbered} after each header except Exerisices
% \renewcommand\thesubsubsection{\arabic{section}}
% \renewcommand\thesubsubsection{\arabic{subsection}}
% \renewcommand\thesubsubsection{\arabic{chapter}-\arabic{subsubsection}}

% scrartcl does not have chapters
% \renewcommand*{\chapterformat}{%
%   \textcolor[rgb]{0.8, 0.8, 0.8}{\thechapter}\autodot\enskip%
% }



% % prevent latex from "floating" the figures
% \usepackage{float}
% \let\origfigure\figure
% \let\endorigfigure\endfigure
% \renewenvironment{figure}[1][2] {
%     \expandafter\origfigure\expandafter[H]
% } {
%     \endorigfigure
% }


%%%%%%% TO DO BEFORE HAND IN %%%%%%%%
%%%%%%%% OR: MAKE KASPER DO %%%%%%%%%

% Add Kasper (supervisor) to front matter
% Make page numbers until TOC roman numerals
% Restart with Latin numbers from page after TOC (Introduction)
% Think about changing fonts for captions, maybe to sans serif
% Looks like Table captions shows two (identical) labels

\usepackage{arxiv}

\usepackage{lmodern}
%\usepackage[hidelinks]{hyperref}       % hyperlinks % clashed with header_extra.tex
\usepackage{url}            % simple URL typesetting
\usepackage{booktabs}       % professional-quality tables
\usepackage{amsfonts}       % blackboard math symbols
\usepackage{nicefrac}       % compact symbols for 1/2, etc.
\usepackage{microtype}      % microtypography
\usepackage{lipsum}		
\usepackage{graphicx}
\usepackage{tikz}
\usepackage[numbers]{natbib}
\usepackage{doi}
\usepackage{amsmath}
\usepackage{amssymb}
\usepackage{svg}
\usepackage{verbatim}
\usepackage{caption}
\usepackage{subcaption}
\usepackage[linesnumbered, ruled, vlined]{algorithm2e}
% Redefine the name for algorithm environment
% \renewcommand{\algorithmautorefname}{Algorithm}
\captionsetup{font=small, labelfont={bf,sf}}
\captionsetup[sub]{font=small, labelfont=bf} 
\captionsetup[subfigure]{justification=justified, singlelinecheck=false, skip=3pt}
\svgsetup{pretex=\tiny}



\title{Chromatin Compartments correlates with Selection}

\author{{\emph{Author:}\\
    Søren Jørgensen}	\\
    \emph{Stud. MSc Bioinformatics} \\
    Bioinformatics Research Center (BiRC),\\
    Dept. of Molecular Biology and Genetics,\\
	Aarhus University, Denmark \\
	\texttt{201906763@post.au.dk} \\
    \And
    {\emph{Supervisor:}\\
    Kasper ~Munch}	\\
    \emph{Associate Professor} \\
    Bioinformatics Research Center (BiRC),\\
    Dept. of Molecular Biology and Genetics,\\
	Aarhus University, Denmark \\
	\texttt{cstorm@birc.au.dk} \\
}

% Uncomment to remove the date
%\date{}

% Uncomment to override  the `A preprint' in the header
\renewcommand{\headeright}{MSc Autumn2024}
\renewcommand{\undertitle}{30 ECTS Master's Thesis, Autumn 2024}
\renewcommand{\shorttitle}{Chromatin Compartments and Selection}
\KOMAoption{captions}{tableheading}
\makeatletter
\@ifpackageloaded{caption}{}{\usepackage{caption}}
\AtBeginDocument{%
\ifdefined\contentsname
  \renewcommand*\contentsname{Table of contents}
\else
  \newcommand\contentsname{Table of contents}
\fi
\ifdefined\listfigurename
  \renewcommand*\listfigurename{List of Figures}
\else
  \newcommand\listfigurename{List of Figures}
\fi
\ifdefined\listtablename
  \renewcommand*\listtablename{List of Tables}
\else
  \newcommand\listtablename{List of Tables}
\fi
\ifdefined\figurename
  \renewcommand*\figurename{Figure}
\else
  \newcommand\figurename{Figure}
\fi
\ifdefined\tablename
  \renewcommand*\tablename{Table}
\else
  \newcommand\tablename{Table}
\fi
}
\@ifpackageloaded{float}{}{\usepackage{float}}
\floatstyle{ruled}
\@ifundefined{c@chapter}{\newfloat{codelisting}{h}{lop}}{\newfloat{codelisting}{h}{lop}[chapter]}
\floatname{codelisting}{Listing}
\newcommand*\listoflistings{\listof{codelisting}{List of Listings}}
\makeatother
\makeatletter
\makeatother
\makeatletter
\@ifpackageloaded{caption}{}{\usepackage{caption}}
\@ifpackageloaded{subcaption}{}{\usepackage{subcaption}}
\makeatother

\usepackage[]{natbib}
\bibliographystyle{plainnat}
\usepackage{bookmark}

\IfFileExists{xurl.sty}{\usepackage{xurl}}{} % add URL line breaks if available
\urlstyle{same} % disable monospaced font for URLs
\hypersetup{
  pdftitle={Chromatin Compartments and Selection on X},
  pdfauthor={Søren Jørgensen},
  pdfkeywords={Hi-C, Chromatin Compartments, Selection},
  hidelinks,
  pdfcreator={LaTeX via pandoc}}


\title{Chromatin Compartments and Selection on X}
\usepackage{etoolbox}
\makeatletter
\providecommand{\subtitle}[1]{% add subtitle to \maketitle
  \apptocmd{\@title}{\par {\large #1 \par}}{}{}
}
\makeatother
\subtitle{How Edges of Active Chromatin Align with Selection Regions in
Primates}
\author{Søren Jørgensen}
\date{2024-11-22}

\begin{document}
\maketitle
\begin{abstract}
This is a dummy abstract, dreamt up by chatGPT. This thesis investigates
the 3D chromatin architecture of the X chromosome in baboons, macaques,
and humans, focusing on chromatin compartments during spermatogenesis.
Using publicly available Hi-C data, interaction maps were created to
identify Principal Component 1 (PC1) compartments, revealing distinct
compartmentalization patterns among species. The analysis included
transition zones, where chromatin shifts between compartment types, and
their correlation with positively selected regions. By comparing these
zones with evolutionarily significant regions, the study explores how
chromatin structure influences evolutionary pressures. Key findings
include conserved chromatin features that may help retain
non-advantageous alleles, suggesting a role for selfish genetic elements
in genome evolution. This research offers new insights into the
relationship between chromatin architecture and evolutionary dynamics
across primate species.
\end{abstract}

\newpage
\tableofcontents

\renewcommand*\contentsname{Table of contents}
{
\setcounter{tocdepth}{2}
\tableofcontents
}

\newpage{}

\section{Introduction}\label{introduction}

{📝} = \emph{draft paragraph}

\subsection{Sexual reproduction (spermatogenesis,
meiosis)}\label{sexual-reproduction-spermatogenesis-meiosis}

{📝} The production of gametes in a sexually reproducing organism is a
highly complex process that involves numeruous elements.
Spermatogenesis, the process of forming male gametes, involves four
stages of differentiation from a germ cell through \emph{spermatogonia},
\emph{pachytene spermatocyte}, and \emph{round spermatids} to
\emph{spermatozoa}, or \emph{sperm} \citep{wang_reprogramming_2019}, and
it is the very basis of male reproduction. The specialized cell division
of meiosis neatly handles the pairing, recombination, and segregation of
homologous chromosomes, thereby ensuring proper genetic distribution.
Deeply understanding the steps of molecular steps of reproduction and
how our genetic material is inherited is essential in biology, bringing
insight to areas such as speciation, population diversity, and (male)
infertility.

\subsection{Selfish genes (and
randomness)}\label{selfish-genes-and-randomness}

{📝} The conventional story of meiosis in gametogenesis is one of random
segregation of the sex chromosomes. They split into haploid gametes,
where each chromosome has an equal chance of being passed on to a
gamete. That seems like a fair game, but what if some genes are cheating
the system by making others less viable. A meiotic driver is a selfish
gene element that modulates meiosis and preferentially transmits its own
allele through meiosis, regardless of the downstream fitness effects it
may have (good or bad) on the organism it is part of. This phenomenon
challenges the traditional understanding of selection, extending its
scope beyond the fitness effects on an organism to include selective
pressures at the molecular level. For example, if some genes on the X
chromosome create a disadvantage for gametes that \emph{do not} contain
those genes, making sure the Y chromosome is not as viable as the X,
resulting in a sex imbalance and possibly numerous other downstream
effects. That is exactly what is coined \emph{sex chromosome meiotic
drive} \citep{jaenike_sex_2001}, a result of selfish genetic elements.

\begin{quote}
{📝} Motivated by previous results in the Munch Research group
\citep{munch_group_2024} on hybrid incompatibility and extended common
haplotypes \citep{skov_extraordinary_2023, sorensen_genome_wide_2023}
that could be explained by meiotic drive, we wanted to investigate how
these patterns correlate with chromatin compartments.
\end{quote}

\subsection{High-Throughput Chromosome Conformation Capture
(Hi-C)}\label{high-throughput-chromosome-conformation-capture-hi-c}

Our DNA can be divided into different orders of structure. \emph{3C}
focus on identifying the highest orders of organization inside the
nucleus, that is, when the 30 nm thick coil of chromatin fibers folds
into loops, Topologically Associating Domains (TADs), and chromatin
compartments. Here, we narrow our focus on the largest of the
structures, \emph{compartments}, that is known to determine availability
to transcription factors, thus making an \emph{A} compartment
\emph{active}---and the \emph{B} compartment \emph{inactive}. The
introduction of the Hi-C (high-throughput 3C) method
\citep{lieberman_aiden_comprehensive_2009} opened new possibilities for
exploring the three-dimensional organization of the genome.

\newpage{}

\section{Methods}\label{methods}

In this project, we formulate two objectives:

\textbf{A}: Reproduce the Hi-C interaction maps and eigendecomposition
from \citep{wang_reprogramming_2019}, with some modifications. We
briefly use \emph{HiCExplorer}, but change the analyses to use the
\emph{Open2C Ecosystem} \citep{open2c} which have a Pyton API as well as
command-line functions, which can be paired very well with Jupyter
Notebooks. The majority of the data analysis was run with a \emph{gwf}
workflow, and the commands that were visually inspected were run in
Jupyter Notebooks.

\textbf{B} Compare with regions of selection that are found in
\emph{papio anubis}, and maybe in \emph{human} too. Investigate the
biological meaning of the results.

The whole of this project is carried out with reproducibility in mind,
so an effort (and quite a significant amount of time) has been put into
documenting code and organizing the project for readbility and
transparency through a Quarto project {[}ref{]}. Therefore, all code and
text is made available as a Quarto book, rendered directly from the
GitHub repository with GitHub Pages {[}{]}. To make this possible, the
Quarto documentation has been extensively studied and discussed with
\emph{KMT}.

\subsection{Downloading Data and Project
Structure}\label{downloading-data-and-project-structure}

To reproduce the results from \citep{wang_reprogramming_2019}, I chose
to use their raw data directly from the SRA portal {[}ref{]}. I filtered
the data to contain all their paired-end Hi-C reads, and included only
macaque samples. The data set also contains RNAseq data, and the same
tissues for both macaque and mouse. The meta data for the data set was
extracted into a runtable \texttt{SRA-runtable.tsv}. To get an overview
of the data accessions used in this analysis, we will first summarize
the runtable that contains the accession numbers and some metadata for
each sample (Table~\ref{tbl-runtable-summary}). It adds up to
\textasciitilde{} 1Tb of compressed \texttt{fastq} files, holding
\textasciitilde{} 9.5 billion reads.

\begin{longtable}[]{@{}lllll@{}}

\caption{\label{tbl-runtable-summary}Summary of the data accessions used
in this analysis}

\tabularnewline

\caption{}\label{T_0aa9c}\tabularnewline
\toprule\noalign{}
~ & source\_name & GB & Bases & Reads \\
\midrule\noalign{}
\endfirsthead
\toprule\noalign{}
~ & source\_name & GB & Bases & Reads \\
\midrule\noalign{}
\endhead
\bottomrule\noalign{}
\endlastfoot
0 & fibroblast & 211.403275 & 553,968,406,500 & 1,846,561,355 \\
1 & pachytene spermatocyte & 274.835160 & 715,656,614,700 &
2,385,522,049 \\
2 & round spermatid & 243.128044 & 655,938,457,200 & 2,186,461,524 \\
3 & sperm & 164.131640 & 428,913,635,400 & 1,429,712,118 \\
4 & spermatogonia & 192.794420 & 518,665,980,300 & 1,728,886,601 \\

\end{longtable}

\subsection{Handling coolers}\label{handling-coolers}

\paragraph{Merge samples from each
source}\label{merge-samples-from-each-source}

\begin{figure}

\centering{

\pandocbounded{\includesvg[keepaspectratio]{index_files/mediabag/svg-xml-base64,PHN2Z.svg}}

}

\caption{\label{fig-flowchart-handling-coolers}A flowchart showing the
pipeline from \texttt{.fastq} to \texttt{.mcool}. Probably BioRender or
Inkscape.}

\end{figure}%

We will use \texttt{cooler\ merge} to merge all samples in each
sub-folder (cell type) to just one interaction matrix for each cell
type. The reason for that is that we choose to trust
\citet{wang_reprogramming_2019} when they say that compartments are
highly reproducible between replicates, and by merging all replicates,
we will have a more robust signal.

Thus, we will merge all samples from the same source\_name into a single
cooler file. The coolers are produced to only contain the real
chromosomes (1-22, X, Y), and not the mitochondrial DNA or the unplaced
contigs (\textasciitilde2,900), as they would not bring any information
to the analysis.

First, we will create a dictionary with the paths to the coolers for
each sample. We use \texttt{glob} to fetch all the coolers in each
sub-folder that remain after filtering and (automatic) quality control
(those are the ones with \texttt{nodups} in their names). \(x\) is the
new \(\frac{x}{2}\)\ldots{}

\newpage{}

\section{Results}\label{results}

Here are the glorious results

\newpage{}

\section{Discussion}\label{discussion}

Here is the discussion

\newpage{}


\renewcommand\refname{Bibliography}
  \bibliography{sub-references.bib}



\end{document}
